\documentclass[11pt]{article}
\usepackage[margin=.84in]{geometry}
\usepackage{amsmath,amssymb,hyperref}
\hypersetup{colorlinks=true,urlcolor=blue}
\newcommand{\Z}{\mathbb Z}
\newcommand{\Q}{\mathbb Q}
\title{\vspace{-1.6em}Exact sparse \(p\)-adic root counts\\ and a sharp exponent-digit bound}
\author{ProveIt GRH article supplement}
\date{27 November 2026}
\begin{document}
\maketitle

The prescribed-support root count admits a finite determinant formula
for \emph{any} number of terms, including more terms than the prime.
A bounded number of base-\(p\) digits of each exponent suffices to
determine the count at every depth. We prove that this digit bound
is sharp for every prime on an infinite family of term counts.
For example, at the prime 2 the first failure of parity alone is
at eight terms: swapping exponent 2 for 8 changes the exact maximum
from seven roots to eight at depth 2.

\section{The determinant rule}
Fix a prime \(p\), distinct nonnegative exponents \(E\) of size \(t\ge2\),
a nonzero center \(A\), and depth \(r>1/(p-1)\).
We allow coefficients and distinct roots over finite extensions of
\(\Q_p\), with \(v_p(p)=1\), in the closed disk \(v_p(X/A-1)\ge r\).
In the basis \(1,T,\ldots,T^{t-1}\), let \(u_k\in\Z^t\) be the remainder
vector of \(T^k\) modulo \(\prod_{n\in E}(T-n)\). Write
\[
 w(k)=rk-v_p(k!),\qquad W=\max_{i<t}w(i),\qquad
 \delta=r-\frac1{p-1}.
\]
Keep indices \(k\le\lfloor W/\delta\rfloor\) with \(w(k)\le W\).
For each \(t\)-element subset \(I\) of those indices, form the cost
\[
 v_p\!\left(\det [u_k]_{k\in I}\right)+\sum_{k\in I}w(k),
\]
with a zero determinant assigned infinite cost.
The exact maximum is the largest index in \emph{any} minimum-cost
set. Thus ties must be retained. The cutoff in the index depends on
\(t,p,r\), not on the largest exponent. The result bounds multiplicities,
and is attainable even by distinct roots with every support coefficient
nonzero. At center \(1\) the coefficients may be chosen rational.

\paragraph{Proof.}
For rational \(r\), choose \(b\) in a finite extension with \(v_p(b)=r\).
Use the \href{https://github.com/VladimirReshetnikov/ProveIt/blob/dd990b27a22290c756c1bd4adce1d971bf7e7821/Analysis/GeneralizedRiemannHypothesis/Article/sections/smd.tex#L57}{article's substitution}
\(X=A\exp(bY)\), identifying the disk with \(v_p(Y)\ge0\).
Its series coefficients are the values of a linear functional on
the scaled columns \(v_k=b^ku_k/k!\). Indeed the unscaled values
on \(u_k\) are the moments \(\sum_{n\in E}c_n A^n n^k\).
After normalizing by a moment of least valuation among the first \(t\),
all moments are integral by the monic recurrence, and some early
series coefficient has valuation at most \(W\). Legendre's formula
gives \(w(k)\ge\delta k\), so discarded indices cannot minimize.
\newpage
Let \(M\) be the lattice generated over the valuation ring by the
retained columns. It is free of rank \(t\).
A subset of columns is a basis exactly when its determinant valuation
is minimum: in any lattice basis every column is integral and some
subset reduces to a basis modulo the maximal ideal, by spanning.
Its determinant is then a unit relative to that basis.
A generating column belongs to some basis exactly when its residue
in \(M/\mathrm m M\) is nonzero.

Normalize a functional by its least valuation on the retained
columns. It is now integral on \(M\), with a unit value at some
column. A column of zero residue cannot give a unit.
Conversely, choose a lattice basis containing the eligible column
of greatest index, and use the corresponding dual functional.
It has a unit value there. Thus the largest possible minimizing index
is the one in the determinant rule. Weierstrass preparation identifies
that index with the root count, including multiplicities.
The determinant comparison is unchanged by any further finite
extension of the coefficient field.

For the dual functional at column \(k\), its coordinates in the
unscaled power basis are rational up to the scalar \(k!/b^k\).
Indeed the underlying rational basis is formed by the \(u_j\)'s,
and the dual row uses only the scaling at \(k\).
At center \(1\) the rational row gives rational polynomial
coefficients by evaluation on the Lagrange cardinal polynomials for
the exponents. Small rational perturbations preserve the minimizing
indices: their contribution to every coefficient has valuation
at least the precision plus \(w(k)\ge0\). One can fill any missing
support coefficients and avoid the discriminant after removing
the smallest monomial. For other centers rescale the variable.
For irrational \(r\) use a single finite candidate set valid in a
neighborhood of \(r\), from the uniform cutoff on a compact interval
above the convergence boundary. The determinant costs are
affine with rational coefficients, so the rational result is constant
in a neighborhood of each irrational depth. Monotonicity then gives
the same count and attainability there.

\section{How many exponent digits suffice?}
Put \(S=\max_{0\le i<t}S_p(i)\), where \(S_p(i)\) is the base-\(p\)
digit sum, and let \(B=\lceil(S-1)/(p-1)\rceil\).
The count at every allowed depth depends only on the multiset of
exponents modulo \(p^B\). To see this, matching that multiset makes
the exponent polynomials congruent modulo \(p^B\), hence makes all
remainder vectors congruent. For \(k\ge t\),
\[
 W-w(k)\le \frac{S-1}{p-1}-\delta < B
\]
by Legendre's formula and \(S_p(k)\ge1\).
Thus the difference between the scaled columns for the two supports
has all coordinate valuations greater than \(W\).
It lies in \(\mathrm m\) times the lattice generated by the first
\(t\) scaled basis vectors, which is identical for the two supports.
Their full lattices coincide and the columns have identical residues
in the common lattice modulo its maximal ideal. This preserves the
answer. The one-term case has no roots in the disk.
\newpage
\section{Sharpness for every prime at \(t=p^m\)}
Take \(t=p^m\), with \(m\ge1\) for odd \(p\) and \(m\ge2\) for \(p=2\).
Here \(S=m(p-1)\), so \(B=m\) for odd \(p\) and \(B=m-1\) for \(p=2\).
Set \(d=1\) for odd \(p\) and \(d=2\) for \(p=2\);
put \(q=p^d\), \(i=t/q=p^{B-1}\).
Compare the two \(t\)-element supports
\[
 E=\{0,\ldots,t-1\},\qquad E'=\{0,\ldots,t\}\setminus\{i\}.
\]
They have equal residue multisets modulo \(p^{B-1}\), since the exchange
\(i\mapsto t\) is zero modulo \(i=p^{B-1}\). Yet at depth \(d\),
and in fact for every \(1/(p-1)<r\le d\), their exact maxima
are \(t-1\) and \(t\), respectively. The first follows by degree
and a product of distinct linear factors close to \(1\) (or by the
default moment witness and small perturbation). For the second, use
\[
 F(X)=(X-1)^{t-1}\bigl((q-1)(X-1)+q\bigr)+p^{dt+1}.
\]
The coefficient of \(X^j\) for \(1\le j<t\) vanishes exactly when
\((q-1)\binom{t}{j}=q\binom{t-1}{j}\), that is, when \(j=i\).
The leading and constant coefficients are nonzero, so the support
is exactly \(E'\).

At depth \(d\), the only coefficient indices of minimum valuation
after the scaling \(X=1+p^dY\) are \(t-1,t\), both with valuation \(dt\).
This puts all \(t\) roots in the disk. The derivative, in \(Z=X-1\),
is \(Z^{t-2}\bigl(t(q-1)Z+q(t-1)\bigr)\). Neither critical point
is a root: \(Z=0\) has nonzero value, and the other point has
valuation \(d-m\le0\). So the roots are distinct.
This proves that the digit bound is optimal for these term counts,
not just an upper bound from the lattice proof.

For \(r>d\), the maximum for \(E'\) drops to \(t-1\).
Otherwise the monic version of a degree-\(t\) polynomial supported
on \(E'\) would have every root \(1+\epsilon\) with \(v_p(\epsilon)>d\).
Its \(X^i\)-coefficient would be congruent to
\((-1)^{t-i}\binom{t}{i}\) up to an error of valuation \(>d\).
But \(\binom{t}{i}=(t/i)\binom{t-1}{i-1}\) has valuation exactly
\(d\), since \((1+X)^{t-1}\bmod p\) has no zero coefficients.
This contradicts the missing term.
The lower bound \(t-1\) follows from the default moment witness.

At \(p=2\), up to seven terms have \(S\le2\), so parity suffices;
the member \(t=8\) of the family is the first counterexample.
At odd \(p\), the family shows that arbitrarily many digits can
be necessary even at the fixed depth \(1\).

\paragraph{Replay and scope.}
The script uses exhaustive determinant costs and cofactor witnesses
on 120 cases, checking the predicted counts by independently shifting
the polynomial and computing Newton weights. It also checks 16 examples
of the sharp digit family. For classical background, see
\href{https://doi.org/10.1112/jlms/s2-36.1.1}{van der Poorten--Rumely}
on zeros of \(p\)-adic exponential polynomials.
This is an unconditional written proof, not a Lean formalization
or a claim of publication priority.
\end{document}
